\section{Experiments}
\label{sec:experiments}

\subsection{Baselines and Training Setup}
We evaluate five representative baselines: AASIST~\cite{Jung2021AASIST}, AASIST-L~\cite{Jung2021AASIST}, XLS-R+SLS~\cite{sls}, SafeEar~\cite{li2024safeear}, and Nes2Net~\cite{Nes2Net}. Together, they cover graph-based anti-spoofing, lightweight anti-spoofing, self-supervised speech representations, content privacy-preserving detection, and recent efficient detector design. Our goal is not to exhaust the detector literature, but to test whether the benchmark reveals delivery-conditioned patterns that are consistent across heterogeneous model families.

We use the official implementations of all baselines and modify only data loading and score export for ChainBench-ADD. AASIST and AASIST-L are trained for 20 epochs with batch size 24 and learning rate $10^{-3}$. XLS-R+SLS, SafeEar, and Nes2Net are trained for 15 epochs using their model-specific batch sizes and learning rate $10^{-6}$. Weight decay is set to $10^{-4}$ for all models.


\subsection{Evaluation Metrics and Settings}
We follow the five benchmark protocols defined in Sec.~\ref{sec:benchmark} and focus here on evaluation setup and metrics.

\noindent\textbf{In-chain detection}
We report family-wise equal error rate (EER, \%) for bona fide-vs.-spoof classification within each delivery family, together with the macro average across families. To assess cross-template generalization, we further adopt a leave-one-template-out protocol for each non-direct family, excluding one template during training and using it only for testing; the reported EER is averaged over held-out templates within that family.

\noindent\textbf{Intervention robustness}
For operator substitution, parameter perturbation, and order swap, we evaluate matched test pairs using the same train/dev pools as in the corresponding benchmark task. We report two primary metrics. The first is pair consistency, i.e., the proportion of pairs for which the model gives the same binary decision to both samples. The second is both-correct rate, i.e., the proportion of pairs for which both samples are classified correctly. Together, these metrics show whether stability under a controlled intervention reflects correct discrimination rather than repeated error. For all decision-based robustness metrics, we use a single EER-defined reference threshold computed once on the pooled evaluation split of each task and keep it fixed for all matched-pair subsets and lineage graphs, without any family- or subgroup-specific retuning. As a secondary score-level diagnostic, we also report the mean normalized score difference (MNSD) between the two samples in each pair. These decision-based metrics should therefore be interpreted as standardized diagnostics at a common operating point rather than as deployment-calibrated results. Higher values are better for pair consistency and both-correct rate, while lower values are better for MNSD. Formal definition details are provided in the appendix.

\noindent\textbf{Path-wise delivery robustness}
We evaluate delivery robustness on lineage graphs. For each lineage, the reference node is chosen as the shortest observed operator-signature realization, and graph distance is measured by shortest path over observed atomic delivery edits. We report the censored first-failure distance (C-FFD), which measures how far a detector remains correct before the first error appears and assigns a maximal value when no failure is observed in the evaluation set. We also compute the fraction of lineages that remain correct within $k$ accumulated delivery edits, denoted by $R_k$, and summarize this robustness curve by its area under the curve (AURC-chain). Larger values indicate more persistent correctness under accumulated delivery edits. Their Formal definitions are also provided in the appendix.


\subsection{Results and Analysis}

\noindent\textbf{Delivery family strongly affects detection difficulty.}
Table~\ref{tab:exp_inchain} reports in-chain detection EER by delivery family. The difficulty of audio deepfake detection varies systematically with the delivery family. Averaged across baselines, direct is the easiest condition (14.80 EER), followed by platform-like (19.12) and telephony (24.81), whereas simulated replay (31.77) and hybrid (31.91) are consistently more difficult. Replay-like recapture and mixed delivery effects introduce more degradation than mild redistribution. Nes2Net and XLS-R+SLS obtain the lowest macro EERs, but all baselines follow the same broad pattern across families. This reveals a stable delivery-conditioned difficulty structure that pooled evaluation would obscure.

\begin{table}[t]
\centering
\footnotesize
\setlength{\tabcolsep}{4pt}
\renewcommand{\arraystretch}{1.02}
\caption{In-chain detection EER (\%) by delivery family. Lower is better.}
\label{tab:exp_inchain}
\vspace{-0.5em}
\begin{tabular}{@{}lcccccc@{}}
\toprule
\textbf{Model} & \textbf{Direct} & \textbf{Platform} & \textbf{Telephony} & \textbf{SimReplay} & \textbf{Hybrid} & \textbf{Macro} \\
\midrule
AASIST       & 13.00 & 24.06 & 27.29 & 39.23 & 42.26 & 29.17 \\
AASIST-L     & 16.89 & 24.49 & 29.30 & 41.57 & 38.79 & 30.21 \\
XLS-R+SLS    & \textbf{12.92} & 11.04 & 16.01 & 20.90 & \textbf{21.11} & 16.40 \\
SafeEar      & 18.14 & 27.43 & 35.65 & 37.93 & 35.94 & 31.02 \\
Nes2Net      & 13.03 & \textbf{8.57} & \textbf{15.82} & \textbf{19.21} & 21.44 & \textbf{15.61} \\
\bottomrule
\end{tabular}
\end{table}

\noindent\textbf{The delivery-family ordering persists under unseen templates.}
Table~\ref{tab:exp_holdout} reports leave-one-template-out in-chain detection. The overall pattern remains close to the standard in-chain results: platform-like is still the easiest family, telephony remains intermediate, and simulated replay/hybrid remain the most difficult. The family ordering therefore appears to reflect effects beyond template memorization and transfers to unseen delivery realizations.

\begin{table}[t]
\centering
\footnotesize
\setlength{\tabcolsep}{4pt}
\renewcommand{\arraystretch}{1.02}
\caption{Leave-one-template-out in-chain detection EER (\%). Lower is better.}
\label{tab:exp_holdout}
\vspace{-0.5em}
\begin{tabular*}{0.96\linewidth}{@{\extracolsep{\fill}}lccccc@{}}
\toprule
\textbf{Model} & \textbf{Platform} & \textbf{Telephony} & \textbf{SimReplay} & \textbf{Hybrid} & \textbf{Macro} \\
\midrule
AASIST       & 23.52 & 28.52 & 41.66 & 41.55 & 33.81 \\
AASIST-L     & 26.53 & 29.04 & 40.68 & 42.27 & 34.63 \\
XLS-R+SLS    & 13.79 & 16.14 & 21.41 & 24.03 & 18.84 \\
SafeEar      & 28.96 & 33.92 & 39.49 & 38.21 & 35.15 \\
Nes2Net      & \textbf{11.33} & \textbf{16.08} & \textbf{21.32} & \textbf{24.09} & \textbf{18.20} \\
\bottomrule
\end{tabular*}
\end{table}

\noindent\textbf{Matched local interventions reveal different failure modes.}
Table~\ref{tab:exp_pair} shifts the focus from family-level variation to controlled local edits. The included three intervention tasks stress detectors in different ways. Operator substitution is the most disruptive intervention overall: averaged across baselines, it yields the lowest both-correct rate and the largest MNSD, indicating that changing operator identity perturbs detector behavior more than changing a parameter within the same operator. Parameter perturbation is the mildest intervention on average, suggesting that many detectors retain partial tolerance to within-operator variation even when they are less robust to larger structural changes in the chain. Order swap exposes a different weakness. It has the highest average pair consistency of the three tasks, yet its average both-correct rate remains well below its consistency level, especially for weaker baselines. The gap between these two quantities is largest for AASIST, AASIST-L, and SafeEar, which suggests that local reordering often preserves a model's decision without preserving its correctness. In other words, these models can be wrong on both elements of a matched pair. By contrast, XLS-R+SLS and Nes2Net show much smaller consistency--correctness gaps across all three interventions, suggesting that their invariance is more often aligned with correct discrimination rather than merely repeated error. Taken together, the matched-paired tasks help distinguish between sensitivity to operator identity, tolerance to moderate parameter variation, and vulnerability to local sequence changes. This failure-mode information is difficult to obtain from the pooled detection metrics alone.

\begin{table}[t]
\centering
\scriptsize
\setlength{\tabcolsep}{3pt}
\renewcommand{\arraystretch}{1.03}
\caption{Matched-pair robustness. Cons.=pair consistency, Both=both-correct rate, and MNSD=mean normalized score difference. Higher is better for Cons./Both; lower is better for MNSD.}
\label{tab:exp_pair}
\vspace{-0.5em}
\begin{tabular}{@{}lccccccccc@{}}
\toprule
\textbf{} & \multicolumn{3}{c}{\textbf{Operator substitution}} & \multicolumn{3}{c}{\textbf{Parameter perturbation}} & \multicolumn{3}{c}{\textbf{Order swap}} \\
\cmidrule(lr){2-4} \cmidrule(lr){5-7} \cmidrule(lr){8-10}
\textbf{Model} & \textbf{Cons.}$\bm{\uparrow}$ & \textbf{Both}$\bm{\uparrow}$ & \textbf{MNSD}$\bm{\downarrow}$ & \textbf{Cons.}$\bm{\uparrow}$ & \textbf{Both}$\bm{\uparrow}$ & \textbf{MNSD}$\bm{\downarrow}$ & \textbf{Cons.}$\bm{\uparrow}$ & \textbf{Both}$\bm{\uparrow}$ & \textbf{MNSD}$\bm{\downarrow}$ \\
\midrule
AASIST     & 74.80 & 54.79 & 52.67 & 78.16 & 62.41 & 43.49 & 80.98 & 58.74 & 41.20 \\
AASIST-L   & 73.75 & 52.97 & 53.10 & 77.22 & 60.88 & 34.66 & 79.19 & 55.56 & 35.16 \\
XLS-R+SLS  & \textbf{87.74} & 81.07 & 20.47 & 90.36 & 85.99 & 16.48 & \textbf{90.48} & \textbf{82.62} & 16.47 \\
SafeEar    & 65.25 & 47.60 & 52.90 & 67.40 & 51.29 & 45.69 & 72.56 & 51.23 & 39.85 \\
Nes2Net    & 87.57 & \textbf{81.86} & \textbf{15.47} & \textbf{90.45} & \textbf{86.92} & \textbf{11.70} & 89.32 & 81.83 & \textbf{13.72} \\
\bottomrule
\end{tabular}
\end{table}

\noindent\textbf{Robustness along delivery paths adds information beyond isolated edits.}
Table~\ref{tab:exp_delivery} evaluates whether correctness persists as delivery edits accumulate along realistic lineage paths. Unlike the matched-pair tasks, which isolate a single local change, this protocol tests stability under multi-step variation. The lineage metrics further separate the two strongest models. Although XLS-R+SLS and Nes2Net are similarly strong under isolated interventions, Nes2Net is consistently better in C-FFD, AURC-chain, and all reported $R_k$ values, suggesting more persistent correctness as delivery edits accumulate. The radius-wise results also show that most newly induced errors appear early: for every model, the largest drop is from $R_0$ to $R_1$, after which the curves flatten substantially, especially for the stronger baselines. Path-wise evaluation therefore shows both when failures first emerge and how quickly they propagate, complementing the view provided by isolated matched-pair tests.

\begin{table}[t]
\centering
\footnotesize
\setlength{\tabcolsep}{4pt}
\renewcommand{\arraystretch}{1.02}
\caption{Lineage-level delivery robustness. C-FFD denotes censored first-failure distance. Higher is better for all metrics.}
\label{tab:exp_delivery}
\vspace{-0.5em}
\begin{tabular*}{0.96\linewidth}{@{\extracolsep{\fill}}lcccccc@{}}
\toprule
\textbf{Model} & \textbf{C-FFD} & \textbf{AURC-chain} & \textbf{$\bm{R_0}$} & \textbf{$\bm{R_1}$} & \textbf{$\bm{R_2}$} & \textbf{$\bm{R_3}$} \\
\midrule
AASIST     & 2.79 & 69.73 & 73.82 & 68.60 & 68.25 & 68.24 \\
AASIST-L   & 2.61 & 65.20 & 69.21 & 64.09 & 63.76 & 63.74 \\
XLS-R+SLS  & 3.47 & 86.76 & 88.39 & 86.31 & 86.18 & 86.18 \\
SafeEar    & 2.46 & 61.45 & 66.97 & 59.93 & 59.45 & 59.44 \\
Nes2Net    & \textbf{3.51} & \textbf{87.86} & \textbf{89.52} & \textbf{87.40} & \textbf{87.27} & \textbf{87.26} \\
\bottomrule
\end{tabular*}
\end{table}

% \noindent\textbf{Summary.}
% The experiments show that ChainBench-ADD captures complementary aspects of delivery-aware audio deepfake detection. Delivery family induces a stable difficulty hierarchy, with direct/platform-like conditions generally easier and simulated replay/hybrid conditions harder. This ordering largely survives leave-one-template-out evaluation, suggesting family-level effects beyond template memorization. The matched intervention tasks reveal distinct local failure modes: operator substitution is typically the most disruptive edit, parameter perturbation is milder, and order swap often separates decision consistency from actual correctness. Lineage-based evaluation adds the cumulative view, showing that many failures emerge within the first one or two edits and that models with similar local robustness can diverge under multi-step delivery variation.
